\documentclass[11pt]{article}
\usepackage[margin=1in]{geometry}
\usepackage{amsmath,amssymb,amsthm}
\usepackage{graphicx}
\usepackage{hyperref}

\newtheorem{theorem}{Theorem}
\newtheorem{lemma}{Lemma}
\newtheorem{corollary}{Corollary}
\newtheorem{proposition}{Proposition}

\newcommand{\R}{\mathbb{R}}
\newcommand{\N}{\mathbb{N}}
\newcommand{\Hol}{\operatorname{Hol}}
\newcommand{\linf}{\ell_\infty}
\newcommand{\norm}[1]{\left\lVert #1\right\rVert}
\newcommand{\abs}[1]{\left\lvert #1\right\rvert}

\title{A One-Dimensional Critical Holder Border Classification for RKHS Sandwiches}
\author{Agent lane 7}
\date{June 15, 2026}

\begin{document}
\maketitle

\section*{Source and Target}

The source is Max Sch\"olpple and Ingo Steinwart,
\emph{Which Spaces can be Embedded in Reproducing Kernel Hilbert Spaces?},
arXiv:2312.14711 \cite{SS}.
In Section 4.2, Theorem 12 proves a packing-number obstruction for RKHS
sandwiches between Holder spaces. For bounded subsets of \(\R^d\) with
non-empty interior, the source concludes that no RKHS can satisfy
\[
  \Hol^\alpha(X)\hookrightarrow H\hookrightarrow \Hol^\beta(X)
\]
when \(\alpha-\beta<d/2\), and no bounded-kernel RKHS can satisfy
\[
  \Hol^\alpha(X)\hookrightarrow H\hookrightarrow \linf(X)
\]
when \(\alpha<d/2\). The paper then says that the bounds are tight up to the
border cases \(\alpha-\beta=d/2\), respectively \(\alpha=d/2\).

\begin{center}
\includegraphics[width=\textwidth]{figures/open_problem_crop.png}
\end{center}

\section*{Partial Result}

This packet gives a one-dimensional classification of the first-order Holder
border cases left open by the source paper. Throughout,
\(\Hol^s([0,1])\) denotes the bounded \(s\)-Holder functions with
norm
\[
  \norm{f}_{s}=\max\left\{\norm{f}_\infty,
  \sup_{x\ne y}\frac{\abs{f(x)-f(y)}}{\abs{x-y}^{s}}\right\},
  \qquad 0<s\le 1 .
\]

\begin{theorem}\label{thm:border}
The one-dimensional critical first-order Holder border has the following
behavior.

\begin{enumerate}
\item Let \(0<\beta<1/2\) and \(\alpha=\beta+1/2\). Then the inclusion
\[
  \Hol^\alpha([0,1])\hookrightarrow \Hol^\beta([0,1])
\]
is not cotype \(2\). Consequently, there is no reproducing kernel Hilbert
space \(H\) on \([0,1]\) such that
\[
  \Hol^\alpha([0,1])\hookrightarrow H\hookrightarrow
  \Hol^\beta([0,1]).
\]

\item The inclusion
\[
  \Hol^{1/2}([0,1])\hookrightarrow \linf([0,1])
\]
is not cotype \(2\). Consequently, there is no bounded-kernel RKHS \(H\) on
\([0,1]\) such that
\[
  \Hol^{1/2}([0,1])\hookrightarrow H\hookrightarrow \linf([0,1]).
\]

\item The endpoint \(\alpha=1\), \(\beta=1/2\) is positive: the Sobolev space
\[
  H^1([0,1])=\{f\in L^2([0,1]): f \text{ has weak derivative }
  f'\in L^2([0,1])\}
\]
with norm
\[
  \norm{f}_{H^1}^2=\norm{f}_{L^2}^2+\norm{f'}_{L^2}^2
\]
is a bounded-kernel RKHS on \([0,1]\), and
\[
  \Hol^1([0,1])\hookrightarrow H^1([0,1])
  \hookrightarrow \Hol^{1/2}([0,1]).
\]
\end{enumerate}
\end{theorem}

Thus the one-dimensional border \(\alpha-\beta=1/2\), \(0\leq\beta\leq1/2\),
has a negative non-Lipschitz side and a positive Lipschitz endpoint. The
remaining higher-dimensional endpoint, not settled here, is the bounded-kernel
case \(\Hol^1([0,1]^2)\hookrightarrow H\hookrightarrow\linf([0,1]^2)\).

\section*{Idea of the Proof}

The source's obstruction tests one metric scale at a time. At the critical
exponent, one scale is exactly balanced: \(2^j\) intervals of length \(2^{-j}\)
and bumps of target norm \(2^{-j/2}\) contribute a bounded amount to the
cotype lower bound. The missing information is that every dyadic scale
contributes this same bounded amount. Taking the first \(m\) scales at once
therefore makes the target \(\ell_2\)-sequence norm grow like \(\sqrt m\).

The domain Rademacher norm stays uniformly bounded because, for
\(\alpha<1\), signed dyadic tents of height \(2^{-j\alpha}\) form uniformly
bounded families in \(\Hol^\alpha\). This forces the cotype-2 constants of the
inclusion to be infinite. Sch\"olpple and Steinwart show that an RKHS sandwich
would make the inclusion 2-factorable, and hence cotype 2; this gives the
contradiction.

The endpoint \(\alpha=1,\beta=1/2\) escapes this obstruction for a real reason,
not just a limitation of the tents. On an interval, every Lipschitz function is
absolutely continuous with essentially bounded derivative, so it belongs to
the Sobolev RKHS \(H^1([0,1])\). Conversely, Cauchy--Schwarz gives
\(H^1([0,1])\hookrightarrow\Hol^{1/2}([0,1])\). This produces an actual
intermediate RKHS at the last one-dimensional border point.

\section*{Dyadic Tent Lemma}

Let
\[
  \theta(t)=\max\{1-\abs{2t-1},0\}, \qquad t\in\R .
\]
Thus \(\theta\) is supported on \([0,1]\), equals \(1\) at \(1/2\), and is
Lipschitz. For \(j\ge 0\) and \(0\le k<2^j\), define
\[
  \theta_{j,k}(x)=2^{-j\alpha}\theta(2^j x-k),\qquad x\in[0,1].
\]

\begin{lemma}\label{lem:tents}
For every \(0<\alpha<1\) there is a constant \(C_\alpha\) such that for every
finite set \(A\subset\{(j,k):j\ge0,\ 0\le k<2^j\}\) and every choice of signs
\(\varepsilon_{j,k}\in\{-1,1\}\),
\[
  \left\|\sum_{(j,k)\in A}\varepsilon_{j,k}\theta_{j,k}
  \right\|_{\Hol^\alpha([0,1])}\le C_\alpha .
\]
\end{lemma}

\begin{proof}
For a fixed level \(j\), set
\[
  g_j(x)=\sum_{k:(j,k)\in A}\varepsilon_{j,k}\theta_{j,k}(x).
\]
The supports at level \(j\) have disjoint interiors and meet only where both
neighboring tents vanish. Hence
\[
  \norm{g_j}_\infty\le 2^{-j\alpha},
  \qquad
  \operatorname{Lip}(g_j)\le C2^{j(1-\alpha)}
\]
with an absolute constant \(C\). The sup norm of \(\sum_j g_j\) is bounded by
\(\sum_j2^{-j\alpha}\).

Let \(x,y\in[0,1]\), \(x\ne y\), and put \(h=\abs{x-y}\). Choose \(J\ge0\) so
that \(2^{-(J+1)}<h\le 2^{-J}\), with the evident modification when
\(h>1/2\). For levels \(j\le J\), the Lipschitz estimate gives
\[
  \sum_{j\le J}\abs{g_j(x)-g_j(y)}
  \le C h\sum_{j\le J}2^{j(1-\alpha)}
  \le C_\alpha h^\alpha .
\]
For levels \(j>J\), the sup-norm estimate gives
\[
  \sum_{j>J}\abs{g_j(x)-g_j(y)}
  \le 2\sum_{j>J}2^{-j\alpha}
  \le C_\alpha h^\alpha .
\]
Combining the two estimates proves the Holder seminorm bound.
\end{proof}

\section*{Proof of the Negative Assertions}

Let \(0<\beta<1/2\) and \(\alpha=\beta+1/2\). For \(m\ge1\), consider the
finite family
\[
  A_m=\{(j,k):0\le j<m,\ 0\le k<2^j\}.
\]
By Lemma \ref{lem:tents},
\[
  \mathbb{E}\left\|
  \sum_{(j,k)\in A_m}\varepsilon_{j,k}\theta_{j,k}
  \right\|_{\Hol^\alpha([0,1])}
  \le C_\alpha ,
\]
where the \(\varepsilon_{j,k}\) are independent Rademacher signs.

On the other hand, each \(\theta_{j,k}\) has a large \(\beta\)-Holder norm.
Indeed, if \(c_{j,k}=(k+1/2)2^{-j}\) is its midpoint and
\(e_{j,k}=k2^{-j}\) is the left endpoint, then
\[
  \frac{\abs{\theta_{j,k}(c_{j,k})-\theta_{j,k}(e_{j,k})}}
       {\abs{c_{j,k}-e_{j,k}}^\beta}
  =
  \frac{2^{-j\alpha}}{(2^{-j-1})^\beta}
  =
  2^\beta 2^{-j(\alpha-\beta)}
  =
  2^\beta 2^{-j/2}.
\]
Therefore
\[
  \left(\sum_{(j,k)\in A_m}
  \norm{\theta_{j,k}}_{\Hol^\beta([0,1])}^2\right)^{1/2}
  \ge
  2^\beta
  \left(\sum_{j=0}^{m-1}2^j2^{-j}\right)^{1/2}
  =
  2^\beta\sqrt m .
\]
The cotype-2 constants of the inclusion
\(\Hol^\alpha([0,1])\hookrightarrow\Hol^\beta([0,1])\) are therefore at least
\(2^\beta C_\alpha^{-1}\sqrt m\) for every \(m\), and hence are unbounded.
The inclusion is not cotype \(2\).

Sch\"olpple and Steinwart prove that if an RKHS \(H\) satisfies
\[
  E\hookrightarrow H\hookrightarrow F,
\]
then the inclusion \(E\hookrightarrow F\) is 2-factorable
\cite[Theorem 6]{SS}. They also record that every 2-factorable operator is
cotype \(2\) \cite[Lemma 5 and Corollary 8]{SS}. Thus the non-cotype result
excludes an intermediate RKHS.

For the bounded-kernel endpoint, take \(\alpha=1/2\) and \(F=\linf([0,1])\).
The same family satisfies
\[
  \norm{\theta_{j,k}}_{\linf([0,1])}=2^{-j/2},
\]
so
\[
  \left(\sum_{(j,k)\in A_m}
  \norm{\theta_{j,k}}_{\linf([0,1])}^2\right)^{1/2}
  =
  \sqrt m .
\]
Lemma \ref{lem:tents} again gives a uniformly bounded Rademacher norm in
\(\Hol^{1/2}([0,1])\). Hence the inclusion
\(\Hol^{1/2}([0,1])\hookrightarrow \linf([0,1])\) is not cotype \(2\), and the
same 2-factorization obstruction excludes every bounded-kernel RKHS sandwich.

\section*{Proof of the Positive Endpoint}

\begin{proposition}\label{prop:h1-endpoint}
The Sobolev space \(H^1([0,1])\), equipped with the norm
\[
  \norm{f}_{H^1}^2=\int_0^1 \abs{f(t)}^2\,dt+
  \int_0^1 \abs{f'(t)}^2\,dt,
\]
is a bounded-kernel RKHS on \([0,1]\). Moreover,
\[
  \Hol^1([0,1])\hookrightarrow H^1([0,1])
  \hookrightarrow \Hol^{1/2}([0,1])
\]
continuously.
\end{proposition}

\begin{proof}
We use the standard continuous representative of each \(H^1\)-class. If
\(f\in\Hol^1([0,1])\), then \(f\) is Lipschitz, hence absolutely continuous,
and \(\abs{f'}\leq \operatorname{Lip}(f)\) almost everywhere. Therefore
\[
  \norm{f}_{H^1}
  \leq \norm{f}_\infty+\operatorname{Lip}(f)
  \leq 2\norm{f}_{\Hol^1([0,1])},
\]
so \(\Hol^1([0,1])\hookrightarrow H^1([0,1])\).

Conversely, let \(f\in H^1([0,1])\). For \(x,y\in[0,1]\), the fundamental
theorem for Sobolev functions and Cauchy--Schwarz give
\[
  \abs{f(x)-f(y)}
  \leq \int_{\min(x,y)}^{\max(x,y)}\abs{f'(t)}\,dt
  \leq \norm{f'}_{L^2}\abs{x-y}^{1/2}.
\]
Also, since the interval has measure \(1\), there is a point \(a\in[0,1]\)
with \(\abs{f(a)}\leq\norm{f}_{L^2}\). Hence
\[
  \abs{f(x)}
  \leq \abs{f(a)}+\int_0^1\abs{f'(t)}\,dt
  \leq \norm{f}_{L^2}+\norm{f'}_{L^2}
  \leq \sqrt2\,\norm{f}_{H^1}.
\]
Thus \(H^1([0,1])\hookrightarrow\Hol^{1/2}([0,1])\) continuously.

The same sup-norm estimate shows that every point evaluation satisfies
\[
  \abs{f(x)}\leq\sqrt2\,\norm{f}_{H^1},\qquad x\in[0,1].
\]
Thus point evaluations are continuous, and \(H^1([0,1])\) is an RKHS. The
uniform bound on the evaluation norms implies that its reproducing kernel is
bounded on the diagonal, hence \(H^1([0,1])\) is a bounded-kernel RKHS.
\end{proof}

\section*{Verification and Novelty Notes}

The proof is deterministic and uses no numerical computation. The key local
estimate is Lemma \ref{lem:tents}; it is included explicitly to avoid relying
on a black-box wavelet characterization of Holder--Besov spaces.

This result is partial relative to the full scope of the source paper, but it
does classify the first-order one-dimensional Holder border. It closes the
non-Lipschitz one-dimensional border negatively, closes the
one-dimensional bounded-kernel endpoint \(\alpha=1/2\) negatively, and closes
the one-dimensional Lipschitz endpoint positively via \(H^1([0,1])\).

This packet does not itself settle the higher-dimensional bounded-kernel
endpoint
\[
  \Hol^1([0,1]^2)\hookrightarrow H\hookrightarrow\linf([0,1]^2),
\]
but a subsequent lane-7 packet on the two-dimensional Lipschitz endpoint
settles that case negatively. The higher-smoothness Besov/Triebel border cases
discussed later in the source paper remain outside this packet.

Before promotion, the run indexes were searched for arXiv id
\texttt{2312.14711}, title keywords, \texttt{RKHS}, \texttt{Holder},
\texttt{H\"older}, \texttt{border cases}, \texttt{alpha=d/2}, and
\texttt{2-factorable}; no existing packet or attempt was found. Bounded web
searches for the source title with \texttt{border cases}, for
\texttt{Holder spaces 2-factorable RKHS}, and for
\texttt{C\string^{1/2} RKHS bounded kernel} found the source arXiv record but
no separate answer to this border case. A later deep pass checked nearby RKHS
regularity and Sobolev-RKHS literature. Those sources gave useful context, but
the endpoint \(H^1\) argument above is elementary and self-contained.

\section*{Human Review Recommendation}

Review should focus on three points. First, check the uniform signed tent
estimate in Lemma \ref{lem:tents}. Second, confirm that the source paper's
``2-factorable implies cotype 2'' obstruction applies to the inclusions with
the same Holder and bounded-kernel conventions used here. Third, verify the
short Sobolev endpoint proof in Proposition \ref{prop:h1-endpoint}; this is
the piece that turns the previous negative packet into a one-dimensional
classification.

\begin{thebibliography}{9}
\bibitem{SS}
Max Sch\"olpple and Ingo Steinwart.
\newblock \emph{Which Spaces can be Embedded in Reproducing Kernel Hilbert Spaces?}
\newblock arXiv:2312.14711, 2023.
\newblock \url{https://arxiv.org/abs/2312.14711}.
\end{thebibliography}

\end{document}
